\section{Native executions and records}\label{sec:model}

We use the frozen native model of \cite{FirstPaper2026}. A machine has $k$
two-way work tapes, alphabet $\Gamma$ containing blank $0$, and $Q$ states.
Thus $Q=\texttt{native.q}+1$, not $\texttt{native.q}$. Put
$a=\max(1,\lceil\log_2|\Gamma|\rceil)$ and
$q=\lceil\log_2 Q\rceil$. Every head starts at $0$ and every work tape is
blank. In a successful transition the machine reads the complete old tuple,
writes at the old work-head positions, moves each head by $-1,0$ or $1$, and
appends at most one bit. A halt test is not a successful transition.

We allow arbitrary integers $d,b\ge1$ and write $P=d^2b$. The canonical-width
length is $L=d^2b(d)$, where $b(d)=\max(1,\lceil\log_2d\rceil)$.
Let $\pi$ send an output ordinal to its demanded raw coordinate; for
transposition it exchanges the row and column indices and preserves the bit
lane. Some replay statements hold for any fixed permutation $\pi$.

For an input $X$, write $x(t),y_h(t),w_h(t,z),s(t)$ for the input position,
work positions, work symbols and state just before successful step $t<T$.
Let $o(t)$ be the number of bits appended strictly before $t$. A return
$(p,t)$ of head $h$ has $p+2\le t<T$, $y_h(p)=y_h(t)$, and no intermediate
visit to that cell. Its excursion is
\[
 \max_{p<u<t}|y_h(u)-y_h(p)|.
\]
Adjacent repeated visits cannot be long returns.

The backward dependency graph has a state edge $t\to t-1$ for $t>0$ and a
work edge from $t$ to each actual last prior visit to its scanned work cell.
A preserving write is still a visit and a version parent. Raw observations
are attached to their read steps; there is no mutable input-tape edge.
Cuts delete whole vertices, including roots and targets. Write
$\operatorname{Read}_X(U)$ for every successful raw read in $U$, and
$\operatorname{Out}_X(V)$ for the actual append steps demanding coordinates
in $V$. Paths of length zero count, so a root that is also a target must be
deleted to separate it.

\begin{lemma}[Old-position writes]\label{lem:old-write}
\formal{NativeWrites.lean}{step_work_update,step_scanned_write,step_preserves_unscanned,instruction_eq_of_old_tuple}
After a successful native step, every unscanned cell is unchanged. At the
scanned cell the new symbol is the instruction's write. Equal old state,
raw symbol and scanned work tuple give the same instruction, regardless of
unobserved tape contents and previously emitted output.
\end{lemma}
\begin{outline}
Unfold the frozen transition definition. It replaces the work symbol exactly
when the queried address equals the old head position. The transition table
does not receive the output history.
\end{outline}

All counts below concern one fixed transition table at each length. Uniform
resource statements allow a nonuniform table only where explicitly stated;
they never take a union over all tables. Correctness means the complete
$P$-bit output, not merely agreement on a selected collection of outputs.
